\section{卷积操作详解}

\subsection{案例研究：检测字母 X}

Amini 使用了一个经典的教学示例来直观解释卷积操作。假设我们要判断一幅图像是否包含字母"X"。即使两个"X"在位移、缩放或旋转后看起来不同，计算机也应该能将它们都识别为"X"。

\begin{figure}[H]
\centering
\includegraphics[width=0.85\textwidth]{slides-images/slide-27.jpg}
\caption{X or X? 图像以像素矩阵表示，但计算机是逐像素比较的——即使是同一个"X"，经过位移后像素值矩阵完全不同。\protect\footnotemark}
\end{figure}
\footnotetext{Slides 第27页。参考 Rohrer "How do CNNs work?"。}

解决方案：不在全局层面比较，而是在\textbf{局部特征}层面比较。一个"X"由若干特征组成——对角线、交叉点等。只要在两幅图像中都能找到这些局部特征，就可以认为它们都是"X"。

\begin{figure}[H]
\centering
\includegraphics[width=0.9\textwidth]{slides-images/slide-28.jpg}
\caption{字母 X 的特征：对角线片段和交叉点。即使两个 X 整体位置不同，但局部特征是匹配的。\protect\footnotemark}
\end{figure}
\footnotetext{Slides 第28页。}

这些局部特征就是\textbf{滤波器}（filters）。每个滤波器是一个小的数字矩阵，代表我们想要检测的某种模式。

\begin{figure}[H]
\centering
\includegraphics[width=0.9\textwidth]{slides-images/slide-29.jpg}
\caption{用于检测"X"特征的三个 3$\times$3 滤波器：分别检测左对角线、十字交叉和右对角线。\protect\footnotemark}
\end{figure}
\footnotetext{Slides 第29页。}

\subsection{卷积运算的数学定义}

卷积操作的步骤如下：

\begin{enumerate}
    \item 将滤波器放置在输入图像的某个位置上
    \item 对滤波器覆盖的区域，执行\textbf{逐元素乘法}（element-wise multiplication）
    \item 将所有乘积\textbf{求和}
    \item 将结果写入\textbf{特征图}（feature map）的对应位置
    \item 将滤波器滑动到下一个位置，重复上述过程
\end{enumerate}

\begin{figure}[H]
\centering
\includegraphics[width=0.9\textwidth]{slides-images/slide-31.jpg}
\caption{卷积操作示例：5$\times$5 输入图像与 3$\times$3 滤波器的卷积。将滤波器在输入上滑动，逐元素相乘后求和，得到特征图。\protect\footnotemark}
\end{figure}
\footnotetext{Slides 第31页。参考 Karn "Intuitive CNNs"。}

以 $5 \times 5$ 的输入图像和 $3 \times 3$ 的滤波器为例，卷积过程生成一个 $3 \times 3$ 的特征图：

\begin{figure}[H]
\centering
\includegraphics[width=0.9\textwidth]{slides-images/slide-32.jpg}
\caption{卷积的第一步：滤波器覆盖左上角 3$\times$3 区域，逐元素相乘后求和得到 4，写入特征图第一个位置。\protect\footnotemark}
\end{figure}
\footnotetext{Slides 第32页。}

对于一般情况，卷积操作可以用如下公式表示：

$$y_{p,q} = \sum_{i=1}^{k} \sum_{j=1}^{k} w_{i,j} \cdot x_{i+p, j+q} + b$$

其中：
\begin{itemize}
    \item $x$ 是输入图像
    \item $w$ 是 $k \times k$ 的滤波器权重矩阵
    \item $b$ 是偏置项
    \item $y_{p,q}$ 是特征图在位置 $(p,q)$ 的输出值
    \item $(p,q)$ 是隐藏层中神经元的位置
\end{itemize}

\begin{figure}[H]
\centering
\includegraphics[width=0.85\textwidth]{slides-images/slide-40.jpg}
\caption{完整的卷积结果：5$\times$5 输入与 3$\times$3 滤波器（stride=1）卷积得到 3$\times$3 的特征图。\protect\footnotemark}
\end{figure}
\footnotetext{Slides 第40页。}

\begin{knowledgebox}{卷积的本质就是"加权求和"}
如果你理解了 Lecture 1 中全连接层的前向传播——对输入进行加权求和再加偏置——那么卷积层做的事情在本质上完全一样。唯一的区别是：全连接层的每个神经元看到\textbf{所有}输入，而卷积层的每个神经元只看到一个\textbf{局部窗口}内的输入，且窗口的权重在空间上共享。
\end{knowledgebox}

\subsection{滤波器产生不同的特征图}

不同的滤波器权重会产生完全不同的特征图。通过改变滤波器的值，可以实现：

\begin{figure}[H]
\centering
\includegraphics[width=0.9\textwidth]{slides-images/slide-41.jpg}
\caption{不同滤波器产生不同效果：锐化（Sharpen）、边缘检测（Edge Detect）和强边缘检测（Strong Edge Detect）。\protect\footnotemark}
\end{figure}
\footnotetext{Slides 第41页。}

\begin{itemize}
    \item \textbf{锐化滤波器}：放大中心像素、减小周围像素的权重，增强图像细节
    \item \textbf{边缘检测滤波器}：本质上是一种导数算子，检测亮度从高到低的突变
    \item \textbf{强边缘检测滤波器}：更激进的导数算子，产生更明显的边缘
\end{itemize}

\begin{knowledgebox}{从手工设计到自动学习滤波器}
在深度学习之前，人们长期使用手工设计的滤波器（如 Sobel 算子、Laplacian 算子等）来检测特定模式。深度学习的关键突破是：\textbf{让滤波器的权重成为可学习的参数}，通过反向传播和梯度下降从数据中自动优化。这样，网络能自动发现对当前任务最有用的特征，而无需人工设计。
\end{knowledgebox}

\subsection{输出尺寸的计算}

卷积操作的输出尺寸取决于三个超参数：

\begin{table}[H]
\centering
\begin{tabular}{lp{10cm}}
\toprule
\textbf{超参数} & \textbf{说明} \\
\midrule
滤波器大小 (kernel size) & 滤波器的空间尺寸，如 $3 \times 3$, $5 \times 5$ \\
步幅 (stride) & 滤波器每次滑动的像素数。stride=1 表示每次移动一个像素，stride=2 表示每次跳两个像素 \\
填充 (padding) & 是否在输入边缘补零。如果不填充，输出尺寸会缩小 \\
\bottomrule
\end{tabular}
\caption{卷积层的关键超参数}
\end{table}

对于输入尺寸为 $n$、滤波器大小为 $k$、步幅为 $s$ 的情况，输出尺寸为：

$$\text{output size} = \left\lfloor \frac{n - k}{s} \right\rfloor + 1$$

\subsection{本章小结}

卷积操作通过滑动滤波器在图像上进行逐元素乘法和求和，生成特征图。不同的滤波器检测不同的模式。在深度学习中，滤波器权重通过训练自动优化，而非手工设计。卷积的三个核心优势是：局部连接、参数共享和保留空间结构。

\newpage
